\documentclass[11pt]{article}

\usepackage[margin=1in]{geometry}
\usepackage{amsmath,amssymb,amsthm}
\usepackage{enumitem}
\usepackage{mathtools}

\setlength{\parindent}{0pt}
\setlength{\parskip}{0.45em}

\title{Homework 2 -- Theory}
\author{}
\date{}

\begin{document}
\maketitle

% ============================================================
\section*{1}

Perceptron begins at \(w=(0,0)\) and predicts \(\operatorname{sign}(w\cdot x)\).
An example \((x,y)\) counts as a mistake when
\[
y\,(w\cdot x)\le 0,
\]
and every mistake updates the weights by \(w\leftarrow w+y\,x\). The two runs
below each make one pass through the data.

The answer to both questions is no, and three points are enough to show it.
Let
\[
x^1=(2,1),\ y^1=+1,
\qquad
x^2=(1,-1),\ y^2=-1,
\qquad
x^3=(0,1),\ y^3=+1 .
\]
These are separable, since \(w^*=(0,1)\) gives \(y^i(w^*\!\cdot x^i)=1\) for
each one.

\paragraph{Order \(S=(x^1,x^2,x^3)\).}

The first example is always a mistake, because \(w\cdot x^1=0\), and the
update gives
\[
w=(0,0)+(2,1)=(2,1).
\]
The second example has
\[
w\cdot x^2=2(1)+1(-1)=1,
\qquad
y^2=-1,
\]
so the prediction has the wrong sign and the weights move again:
\[
w=(2,1)-(1,-1)=(1,2).
\]
The third example has
\[
w\cdot x^3=1(0)+2(1)=2,
\qquad
y^3=+1,
\]
which is correct, so nothing changes. This run ends with
\[
\boxed{
S:\quad 2\text{ mistakes},\qquad w=(1,2).
}
\]

\paragraph{Order \(S'=(x^3,x^1,x^2)\).}

The first example is again a mistake, since \(w\cdot x^3=0\), and now
\[
w=(0,0)+(0,1)=(0,1).
\]
The second example has
\[
w\cdot x^1=0(2)+1(1)=1,
\qquad
y^1=+1,
\]
which is correct, and the third has
\[
w\cdot x^2=0(1)+1(-1)=-1,
\qquad
y^2=-1,
\]
which is also correct. This run ends with
\[
\boxed{
S':\quad 1\text{ mistake},\qquad w=(0,1).
}
\]

The two orders give different numbers of mistakes and different final
weights, and since \((1,2)\) is not a multiple of \((0,1)\), they describe
different halfspaces as well. Both vectors already classify all three points
correctly, so running until convergence would not change either result.

% ============================================================
\section*{2}

Let \(z=y\,w^Tx\), so that \(z>0\) means the example is classified correctly
and \(z<0\) means it is a mistake. The loss
\[
\phi(z)=\max(0,-z)
=
\begin{cases}
0, & z\ge0,\\
-z, & z<0,
\end{cases}
\]
is therefore zero on correct examples and equal to \(|z|\) on mistakes, and
its derivative is
\[
\phi'(z)
=
\begin{cases}
0, & z>0,\\
-1, & z<0.
\end{cases}
\]
The problem says not to worry about \(z=0\), so treat it as a mistake and set
\(\phi'(0)=-1\). Perceptron does the same thing, updating on its first
example even though \(w\cdot x=0\) there.

Since \(\nabla_w(y^iw^Tx^i)=y^ix^i\), the chain rule gives
\[
\nabla_w\,\phi(y^iw^Tx^i)
=
\begin{cases}
0, & y^iw^Tx^i>0,\\
-y^ix^i, & y^iw^Tx^i<0,
\end{cases}
\]
and substituting this into \(w_{\mathrm{new}}=w_{\mathrm{old}}-\eta\nabla_w\phi\)
leaves
\[
w_{\mathrm{new}}
=
\begin{cases}
w_{\mathrm{old}}, & \text{correct},\\
w_{\mathrm{old}}+\eta\,y^ix^i, & \text{mistake}.
\end{cases}
\]
Setting \(\eta=1\) turns this into
\[
\boxed{
w_{\mathrm{new}}
=
\begin{cases}
w_{\mathrm{old}}, & \text{correct},\\
w_{\mathrm{old}}+y^ix^i, & \text{mistake},
\end{cases}
}
\]
which is the Perceptron update rule.

The step size does not matter either. Every update is multiplied by \(\eta\)
and both algorithms start from the same weights, so the SGD weights remain
\(\eta\) times the Perceptron weights throughout. Scaling by a positive
constant does not change \(\operatorname{sign}(w\cdot x)\), so the two
algorithms predict the same labels and make the same mistakes.

Perceptron is therefore SGD on \(\frac1m\sum_i\phi(y^iw^Tx^i)\), started at
\(w=0\) with \(\eta=1\), taking the examples in order rather than at random.

% ============================================================
\section*{3}

\subsection*{(a)}

The target \(c=h_{\theta_1,\theta_2,b}\) divides the line into three pieces,
\[
L=(-\infty,\theta_1),
\qquad
M=[\theta_1,\theta_2],
\qquad
R=(\theta_2,\infty),
\]
where \(c\) equals \(-b\) on \(L\), \(b\) on \(M\), and \(-b\) again on
\(R\). Writing \(p_L,p_M,p_R\) for the probabilities of the three pieces,
\[
p_L+p_M+p_R=1 .
\]

A stump has a single threshold, so it splits the line in two. It can match
\(c\) on two of the pieces but has to be wrong on the third, which leaves
three choices:
\begin{itemize}[leftmargin=*,itemsep=2pt]
    \item \(h_{\theta_2,\,b}\) predicts \(b\) on \(L\cup M\) and \(-b\) on
    \(R\), so it is wrong only on \(L\) and has error \(p_L\).
    \item \(h_{\theta_1,\,-b}\) predicts \(-b\) on \(L\) and \(b\) on
    \(M\cup R\), so it is wrong only on \(R\) and has error \(p_R\).
    \item \(h_{B,\,-b}\) predicts \(-b\) across all of \([-B,B]\), so it is
    wrong only on \(M\) and has error \(p_M\).
\end{itemize}

Three numbers that sum to \(1\) cannot all be larger than \(\frac13\), so
\[
\min(p_L,p_M,p_R)\le\frac13 ,
\]
and taking the stump that matches the smallest of them gives
\[
\boxed{
\text{some } h\in H \text{ satisfies }
\Pr[h(x)\neq c(x)]\le\frac13 .
}
\]

\subsection*{(b)}

Sort the sample so that \(x_{(1)}\le\cdots\le x_{(m)}\), with labels
\(y_{(1)},\ldots,y_{(m)}\). A stump predicts \(b'\) on the first \(k\) sorted
points and \(-b'\) on the rest, and any threshold between \(x_{(k)}\) and
\(x_{(k+1)}\) gives the same predictions. Only the pair \((k,b')\) matters,
which leaves \(2(m+1)\) candidates, each with error
\[
\mathrm{err}_S(k,b')
=
\#\{j\le k:y_{(j)}=-b'\}
+
\#\{j>k:y_{(j)}=b'\} .
\]

A single sweep finds all of these errors. At \(k=0\) every point is predicted
\(-b'\), so \(\mathrm{err}_S(0,b')=\#\{j:y_{(j)}=b'\}\), and raising \(k\) by
one moves \(x_{(k)}\) to the side predicted \(b'\), which fixes that point if
its label is \(b'\) and breaks it otherwise:
\[
\mathrm{err}_S(k,b')
=
\mathrm{err}_S(k-1,b')
+
\begin{cases}
-1, & y_{(k)}=b',\\
+1, & y_{(k)}=-b'.
\end{cases}
\]

The procedure is to sort the sample, sweep \(k\) from \(0\) to \(m\) for
\(b'=+1\) and \(b'=-1\) while keeping the smallest error, and return the
corresponding stump. Sorting costs \(O(m\log m)\) and each sweep costs
\(O(m)\), for a total of
\[
\boxed{
O(m\log m) .
}
\]
The result has the smallest training error of any stump, so it is an ERM.

\subsection*{(c)}

A stump can label \(m\) points in only \(2(m+1)\) ways, so there are very few
hypotheses to consider, and the training error of each one stays close to its
true error once \(m\) is reasonably large. The stump from part (b) therefore
has true error close to \(\frac13\).

The reason this is easier than learning \(C\) is the accuracy being asked
for. Learning \(C\) would mean driving the error down to \(\epsilon\) and
locating both \(\theta_1\) and \(\theta_2\), so the sample size would have to
grow as \(\epsilon\) shrinks. A weak learner for \(H\) only needs a fixed
advantage over \(\frac12\), which a fixed sample size already gives. Boosting
can then take that advantage and drive the error down to \(\epsilon\), so
\(C\) is weakly learnable using \(H\).

% ============================================================
\section*{4}

At iteration \(t\), each example \(x_i\) carries a weight \(w_i\), the total
weight is \(W=\sum_i w_i\), and the distribution is \(D_t(i)=w_i/W\). The
hypothesis \(h_t\) returned at this iteration has
\[
E_t=\frac1W\sum_{i:\,h_t(x_i)\neq y_i}w_i,
\qquad
A_t=1-E_t,
\qquad
\beta_t=\frac{E_t}{A_t},
\]
and the reweighting multiplies the weight of each correct example by
\(\beta_t\) while leaving the wrong ones as they are:
\[
w_i^{\mathrm{new}}
=
\begin{cases}
\beta_t\,w_i, & h_t(x_i)=y_i,\\
w_i, & h_t(x_i)\neq y_i .
\end{cases}
\]
Assume \(0<E_t<1\), so that \(\beta_t\) is defined and positive.

Before the update the correct examples hold \(A_tW\) of the weight and the
wrong ones hold \(E_tW\). Scaling the correct side gives
\[
\beta_t\cdot A_tW
=
\frac{E_t}{A_t}\cdot A_tW
=
E_tW,
\]
which is exactly what the wrong side already holds, so the new total is
\[
W^{\mathrm{new}}=E_tW+E_tW=2E_tW,
\]
and the accuracy of \(h_t\) on the new distribution is
\[
\boxed{
\Pr_{i\sim D_{t+1}}[h_t(x_i)=y_i]
=
\frac{E_tW}{2E_tW}
=
\frac12 .
}
\]

The factor \(\beta_t\) is chosen precisely so that the correct side shrinks
until it matches the wrong side, and both \(E_t\) and \(W\) cancel, so the
answer is exactly \(\frac12\) for any \(h_t\) and any set of weights. This
leaves \(h_t\) no better than guessing on \(D_{t+1}\), so the next weak
learner has to get a different set of examples right.

\end{document}